%=============================================================================
% Section 5: AI Protocol and Capability Definition
%=============================================================================
\section{AI Protocol and Capability Definition}

\subsection{Protocol Role}

The AI Protocol defines how capabilities are described. It specifies:
\begin{itemize}
    \item \textbf{Providers:} Entities offering AI or non-AI capabilities
    \item \textbf{Models:} Specific AI models with defined capabilities
    \item \textbf{Capability Types:} Categories of capabilities (chat, embedding, image generation, etc.)
    \item \textbf{Endpoints:} Where capabilities can be invoked
    \item \textbf{Parameters:} Input and output schemas
\end{itemize}

The protocol serves as the contract boundary between orchestration and execution: orchestration decides policy, while execution enforces invocation semantics. The protocol is typically represented as structured manifests:

\begin{verbatim}
provider/
    +-- metadata.yaml
    +-- models/
    |     +-- model_a.yaml
    |     +-- model_b.yaml
    +-- capabilities/
          +-- chat.yaml
          +-- embedding.yaml
\end{verbatim}

\subsection{Capability Definition}

Each capability is defined with:
\begin{itemize}
    \item \textbf{Name:} Semantic identifier (e.g., \texttt{generate\_text}, \texttt{search\_web})
    \item \textbf{Input Schema:} Structured input specification
    \item \textbf{Output Schema:} Structured output specification
    \item \textbf{Execution Semantics:} Invocation and runtime behavior contract
    \item \textbf{Quality Metadata ($Q$):} Measured quality fields emitted by execution
    \item \textbf{Pricing Metadata ($P$):} Metering and billing fields emitted by execution
    \item \textbf{Provider Metadata:} Information about the provider offering this capability
\end{itemize}

\textbf{Example Capability Definition:}

\begin{lstlisting}[style=yamlstyle]
name: generate_text
description: Generate text based on a prompt
input_schema:
  type: object
  properties:
    prompt:
      type: string
      description: The input prompt
    max_tokens:
      type: integer
      description: Maximum tokens to generate
  required: [prompt]
output_schema:
  type: object
  properties:
    text:
      type: string
      description: Generated text
execution:
  endpoint: /v1/completions
  method: POST
quality:
  telemetry_fields: [execution_latency_ms, retry_count]
pricing:
  billing_model: per_token
  currency: USD
provider:
  name: Example Provider
  version: 1.0.0
\end{lstlisting}

\subsection{Provider-Centric Organization}

An important design decision is that capability descriptions remain provider-centric rather than capability-centric. The manifest structure follows:

\begin{verbatim}
Provider -> Models -> Capabilities
\end{verbatim}

Rather than:

\begin{verbatim}
Capability -> Providers
\end{verbatim}

The runtime can dynamically construct a capability index:

\begin{verbatim}
chat/
    +-- provider A
    +-- provider B
    +-- provider C
embedding/
    +-- provider A
    +-- provider D
\end{verbatim}

This index can be rebuilt periodically and optimized based on:
\begin{itemize}
    \item User preferences
    \item Latency measurements
    \item Cost calculations
    \item Reliability metrics
\end{itemize}
